\documentclass[../../main.tex]{subfiles}
\graphicspath{{../../image/}{../../figure/}}
\begin{document}

\section{真正非线性性与线性退化}\label{section:ch07-真正非线性}

区分线性与非线性双曲守恒律组行为的特征在于：在前者中，特征速度为常数，同一族的所有波以固定速度传播；而在后者中，波速随波幅变化。随着研究的推进，我们将遇到非线性性的各种表现，并且在每一情形中都会注意到：当特征速度 $\lambda_i$ 沿相应特征向量 $R_i$ 的方向变化时，其效应将格外显著。这激发了下述定义。

\begin{definition}[真正非线性性与线性退化]\label{definition:ch07-真正非线性}
对于 $\mathscr{O}$ 上的双曲守恒律组\eqref{align:ch07-典范守恒律}，$U\in\mathscr{O}$ 称为 $i$-特征族的\emph{真正非线性态}，如果
\begin{align}
\mathrm{D}\lambda_i(U)R_i(U)\neq 0,
\label{align:ch07-真正非线性条件}
\end{align}
或称为 $i$-特征族的\emph{线性退化态}，如果
\begin{align}
\mathrm{D}\lambda_i(U)R_i(U)=0.
\label{align:ch07-线性退化条件}
\end{align}
当对所有 $U\in\mathscr{O}$ 都成立\eqref{align:ch07-真正非线性条件}时，$i$ 是\emph{真正非线性特征族}；而当对所有 $U\in\mathscr{O}$ 都满足\eqref{align:ch07-线性退化条件}时，$i$ 是\emph{线性退化特征族}。当每一特征族都是真正非线性的，\eqref{align:ch07-典范守恒律}是\emph{真正非线性方程组}。
\end{definition}

显然，$i$-特征族是线性退化的，当且仅当 $i$-特征速度 $\lambda_i$ 沿向量场 $R_i$ 的积分曲线为常数。

标量守恒律\eqref{align:ch07-标量守恒律}（特征速度 $\lambda=f'(u)$）当 $f$ 无拐点时是真正非线性的：$f''(u)\neq 0$。特别地，Burgers 方程\eqref{align:ch04-Burgers}
是真正非线性的。

利用\eqref{align:ch07-p系统特征值}与\eqref{align:ch07-p系统右特征}，可直接验证：当 $\sigma''(u)\neq 0$ 时，方程组\eqref{align:ch07-p系统}是真正非线性的。作为习题，读者可验证：当 $2p'(\rho)+\rho p''(\rho)>0$ 时，方程组\eqref{align:ch07-等熵Euler}是真正非线性的，于是特别地，理想气体的方程组\eqref{align:ch07-理想气体}是真正非线性的。浅水波方程组\eqref{align:ch07-浅水波}同样是真正非线性的。

由\eqref{align:ch07-热弹性特征速度}，热弹性方程组\eqref{align:ch07-热弹性}的 2-特征族是线性退化的。结果表明，只要 $\sigma_{uu}(u,s)\neq 0$，另外两个特征族是真正非线性的。

接下来考虑平面弹性振荡方程组\eqref{align:ch07-平面振荡}（沿方向 $\nu$），回忆 $\sigma(u)=\partial e(u)/\partial u$，$e(u)$ 凸。六个特征速度是 $e(u)$ 的 Hessian 矩阵（即在 $F=I+u\nu^\top$ 处赋值的声学张量\eqref{align:ch03-声张量}
的特征值）$\mu_1(u),\mu_2(u),\mu_3(u)$ 的平方根 $\pm\sqrt{\mu_1},\pm\sqrt{\mu_2},\pm\sqrt{\mu_3}$。简单计算表明，关联于特征速度 $\pm\sqrt{\mu_\ell}$ 的特征族在 $u=F\nu$ 处是真正非线性的，如果
\begin{align}
\sum_{i,j,k=1}^3\frac{\partial^3 e(u)}{\partial u_i\partial u_j\partial u_k}\xi_i\xi_j\xi_k
=\sum_{i,j,k=1}^3\sum_{\alpha,\beta,\gamma=1}^3\frac{\partial^3\varepsilon(F)}{\partial F_{i\alpha}\partial F_{j\beta}\partial F_{k\gamma}}\xi_i\xi_j\xi_k\nu_\alpha\nu_\beta\nu_\gamma\neq 0,
\label{align:ch07-弹性三阶导}
\end{align}
其中 $\xi$ 是声学张量关联于特征值 $\mu_\ell$ 的特征向量。

把上述应用于特殊方程组\eqref{align:ch07-各向同性不可压缩}，发现 $\mu_1=h''(|u|)$ 是单特征值（对应特征向量 $u$），而 $\mu_2=\mu_3=h'(|u|)/|u|$ 是二重特征值，其特征空间是 $u$ 的正交补。于是，特征速度 $\pm[h''(|u|)]^{1/2}$ 关联于纵向振荡，而 $\pm[h'(|u|)/|u|]^{1/2}$ 关联于横向振荡。然而，只有同时正交于 $\nu$ 的横向振荡才与不可压缩性相容。当 $h'''(|u|)\neq 0$ 时，关联于 $\pm[h''(|u|)]^{1/2}$ 的特征族在 $u$ 处是真正非线性的；而关联于 $\pm[h'(|u|)/|u|]^{1/2}$ 的特征族是线性退化的。显然，同样的结论适用于弹性弦振荡方程组\eqref{align:ch07-弹性弦}，只需把 $h'(|u|)$ 换成 $\tau(|u|)$。对该方程组，所有横向振荡在物理上都是有意义的，因为不可压缩约束不再相关。模型方程组\eqref{align:ch07-脐点模型}表现出类似行为：它的 1-特征族是线性退化的，而 2-特征族是真正非线性的，原点除外。

在三相流通过多孔介质的方程组\eqref{align:ch07-三相流}中，真正非线性性沿三条对称线 $u=v$、$u=1-u-v$ 与 $v=1-u-v$，以及沿包围点 $u=v=1/3$ 的闭曲线失效。

最后，在电泳方程组\eqref{align:ch07-电泳}中，$n$-特征族是线性退化的，而其余族是真正非线性的。

Born--Infeld 介质\eqref{align:ch03-BornInfeld}
的 Maxwell 方程\eqref{align:ch03-Maxwell}
方程组具有显著性质：在任一空间方向 $\nu\in\mathbb{S}^2$ 上的平面振荡由一个特征族全为线性退化的方程组控制。

相当常见地，线性退化源于严格双曲性的丧失。事实上，\eqref{align:ch07-退化点Dlambda}的一个直接推论是：

\begin{theorem}\label{theorem:ch07-多重特征线性退化}
在双曲守恒律组\eqref{align:ch07-典范守恒律}中，假定 $j$- 与 $k$-特征速度重合：$\lambda_j(U)=\lambda_k(U)$，$U\in\mathscr{O}$。那么 $j$- 与 $k$-特征族都是线性退化的。
\end{theorem}

当方程组\eqref{align:ch07-典范守恒律}配有 Riemann 不变量坐标系 $(w_1,\cdots,w_n)$ 并采用 $W=(w_1,\cdots,w_n)^\top$ 作为状态向量时，真正非线性性与线性退化的条件呈现优雅而富有启发性的形式。事实上，利用\eqref{align:ch07-链式法则}，我们推出\eqref{align:ch07-真正非线性条件}与\eqref{align:ch07-线性退化条件}分别等价于
\begin{align}
\frac{\partial\lambda_i}{\partial w_i}\neq 0
\label{align:ch07-真正非线性W}
\end{align}
与
\begin{align}
\frac{\partial\lambda_i}{\partial w_i}=0.
\label{align:ch07-线性退化W}
\end{align}

\end{document}
